\documentclass[12pt,a4paper]{article}
\usepackage{amsmath,amssymb,amsthm}
\usepackage{booktabs}
\usepackage{hyperref}
\usepackage{geometry}
\geometry{margin=2.5cm}
\usepackage{enumitem}

\newtheorem{theorem}{Theorem}[section]
\newtheorem{lemma}[theorem]{Lemma}
\newtheorem{corollary}[theorem]{Corollary}
\newtheorem{proposition}[theorem]{Proposition}
\theoremstyle{definition}
\newtheorem{definition}[theorem]{Definition}
\theoremstyle{remark}
\newtheorem{remark}[theorem]{Remark}

\title{The Origin of CP Violation from Recognition Science:\\
CPT Invariance, Ledger Phase Structure, and the 8-Tick Epoch}

\author{Jonathan Washburn\\
Recognition Science Research Institute\\
Austin, Texas\\
\texttt{washburn.jonathan@gmail.com}}

\date{March 2026}

\begin{document}
\maketitle

\begin{abstract}
We derive the origin of $CP$ violation from the Recognition
Science (RS) ledger structure~\cite{RS_Axioms}.  $CPT$ invariance is automatic:
the combined operation of charge conjugation ($C$: charge $\to$
$-$charge), parity ($P$: position $\to$ $-$position), and time
reversal ($T$: tick $\to (7 - \mathrm{tick}) \bmod 8$) preserves
the cost functional $J(x)$ and ledger balance.  $CPT$ is
involutive ($\mathrm{CPT}^2 = \mathrm{id}$).  $CP$ violation
arises because $C$ and $P$ individually preserve cost but their
product $CP$ does not preserve the 8-tick phase ordering: the
forward and backward phase sequences are inequivalent when
weak-sector entries carry complex phases from the three-generation
CKM matrix.  Since $CPT$ is exact, $CP$ violation is the
complement of $T$ violation.  The $CPT$ mass bound
$|\Delta m/m| < 10^{-9}$ is satisfied automatically.  We prove
that $CP$ violation requires $N_{\mathrm{gen}} \geq 3$ (two
generations give no $CP$-odd phase), connecting this result to the
RS prediction $N_{\mathrm{gen}} = 3$.  The \emph{existence} of $CP$
violation is a structural consequence; the \emph{magnitude}
$\mathcal{J} \approx 3.2 \times 10^{-5}$ is an open calculation.
\end{abstract}

%% ============================================================
\section{Introduction}
\label{sec:intro}
%% ============================================================

$CP$ violation---matter--antimatter asymmetry in weak interactions---was
discovered in 1964 in neutral kaon decays~\cite{Christenson1964}.
The Standard Model accommodates $CP$ violation through the
Kobayashi--Maskawa mechanism: with $N_{\rm gen} \geq 3$, the quark
mixing matrix carries an irreducible complex phase~\cite{KM1973}.

However, the Standard Model does not explain:
\begin{itemize}
  \item Why $CP$ is violated (why the phase is non-zero).
  \item Why $CPT$ is preserved (the $CPT$ theorem relies on
  Lorentz invariance, locality, and unitarity---assumptions
  about the continuum).
  \item Why there are exactly three generations.
\end{itemize}

RS answers all three from the ledger structure.

%% ============================================================
\section{Recognition Science Framework}
\label{sec:framework}
%% ============================================================

RS derives all physical law from three axioms~\cite{RS_Axioms}:
\begin{enumerate}[label=(A\arabic*),nosep]
  \item $J(1) = 0$ \hfill\textit{(Normalization)}
  \item $J(xy) + J(x/y) = 2J(x)J(y) + 2J(x) + 2J(y)$ for all $x,y>0$
  \hfill\textit{(Recognition Composition Law)}
  \item $J''_{\log}(0) = 1$, where $J_{\log}(t) := J(e^t)$
  \hfill\textit{(Calibration)}
\end{enumerate}

\begin{theorem}[Cost Uniqueness~\cite{RS_Axioms}]
\label{thm:J}
The unique continuous solution is $J(x) = \tfrac{1}{2}(x+x^{-1})-1$,
satisfying $J(x) \geq 0$ (equality iff $x=1$) and $J(x) = J(x^{-1})$.
\end{theorem}

\begin{proof}[Proof sketch]
$H(t) := 1+J(e^t)$ satisfies the d'Alembert equation $H(s+t)+H(s-t)=2H(s)H(t)$
with $H(0)=1$, $H''(0)=1$.  By the Acz\'el classification, $H(t)=\cosh t$,
giving $J(x) = \tfrac{1}{2}(x+x^{-1})-1$.
\end{proof}

Each ledger entry $e$ carries a positive ratio $x_e > 0$ (the recognition
scale ratio between debit and credit).  The discrete symmetries act on
these ratios and on the phase structure of the 8-tick epoch.

%% ============================================================
\section{Discrete Symmetries on the Ledger}
\label{sec:symmetries}
%% ============================================================

\begin{definition}[Charge Conjugation $C$]
On the recognition ledger, $C$ maps charge $q \to -q$
(equivalently, it swaps debit and credit entries).  Since
$J(x)$ depends on $|q|$ through the gap weight, $J$ is
invariant: $J(C(e)) = J(e)$.
\end{definition}

\begin{definition}[Parity $P$]
$P$ maps spatial position $\mathbf{x} \to -\mathbf{x}$.  The
cost functional depends on $|\mathbf{x}|$ (distance), not
direction, so $J(P(e)) = J(e)$.
\end{definition}

\begin{definition}[Time Reversal $T$]
On the 8-tick epoch, $T$ maps tick $t \to (7 - t) \bmod 8$.
This reverses the phase sequence within the epoch.
$T$ preserves cost: $J(T(e)) = J(e)$.
\end{definition}

%% ============================================================
\section{$CPT$ Invariance}
\label{sec:cpt}
%% ============================================================

\begin{theorem}[$CPT$ Preservation]
\label{thm:cpt}
$CPT$ preserves cost:
\begin{equation}
  J(\mathrm{CPT}(e)) = J(e)
\end{equation}
for every ledger entry $e$.  $CPT$ also preserves ledger balance
(total charge conservation).
\end{theorem}

\begin{proof}
We verify cost-preservation for each transformation individually.

\textit{Charge conjugation $C$}: $C$ negates the charge $q \to -q$,
which at the ratio level maps $x_e \to x_e^{-1}$ (debits become credits
and vice versa).  By Theorem~\ref{thm:J}, $J(x^{-1}) = J(x)$, so
$J(C(e)) = J(x_e^{-1}) = J(x_e) = J(e)$.

\textit{Parity $P$}: $P$ maps $\mathbf{x} \to -\mathbf{x}$.  The cost
functional $J(x_e)$ depends on the scale ratio $x_e > 0$, not on the
spatial position.  Hence $J(P(e)) = J(e)$.

\textit{Time reversal $T$}: $T$ maps tick $t \to (7-t)\bmod 8$, reversing
the epoch phase ordering.  $J(x_e)$ depends on $x_e$, not on the tick
index.  Hence $J(T(e)) = J(e)$.

Composition: since all three preserve $J$, so does their product:
$J(\mathrm{CPT}(e)) = J(e)$.

Ledger balance is preserved because $C$ swaps debits and credits
(preserving the total), $P$ and $T$ do not alter charge values.
\end{proof}

\begin{theorem}[$CPT$ Involutive]
\label{thm:involutive}
$\mathrm{CPT}^2 = \mathrm{id}$.
\end{theorem}

\begin{proof}
$C^2 = \mathrm{id}$ (charge double-negation: $q \to -q \to q$).
$P^2 = \mathrm{id}$ (position double-negation:
$\mathbf{x} \to -\mathbf{x} \to \mathbf{x}$).
$T^2 = \mathrm{id}$ (tick reversal:
$t \to (7-t) \to (7-(7-t)) = t \bmod 8$).
Since $C$, $P$, $T$ act on independent degrees of freedom (charge, position, and tick, respectively), they commute: $CPT = CTP = PCT = \ldots$.  Therefore
$(\mathrm{CPT})^2 = C^2 P^2 T^2 = \mathrm{id}$.
\end{proof}

\begin{corollary}[$CPT$ Consequences]
$CPT$ invariance implies:
\begin{enumerate}[label=(\roman*)]
  \item $m_{\mathrm{particle}} = m_{\mathrm{antiparticle}}$.
  \item $\tau_{\mathrm{particle}} = \tau_{\mathrm{antiparticle}}$
  (equal lifetimes).
  \item $q_{\mathrm{particle}} = -q_{\mathrm{antiparticle}}$
  (opposite charges).
\end{enumerate}
\end{corollary}

%% ============================================================
\section{Why $CP$ Is Violated}
\label{sec:cp-violation}
%% ============================================================

\begin{proposition}[$CP$ Violation from Phase Structure]
\label{prop:cp}
When weak-sector ledger entries carry complex CKM phases (imported from
the Standard Model), $CP$ does not preserve the 8-tick phase ordering:
the forward phase sequence and the $CP$-transformed sequence are
inequivalent whenever $N_{\mathrm{gen}} \geq 3$.
\end{proposition}

\begin{proof}
The 8-tick epoch has a natural forward direction: ticks proceed
$0, 1, 2, \ldots, 7$ with phases $e^{2\pi i k/8}$ for
$k = 0, \ldots, 7$.  Under $CP$, the charge and spatial parity
are reversed, but the tick ordering is preserved.  The weak-sector
entries carry generation-dependent phases from the CKM matrix:
$V_{ij} = |V_{ij}| e^{i\delta_{ij}}$.  Under $CP$:
$V_{ij} \to V_{ij}^*$ (complex conjugation).  If
$\mathrm{Im}(V_{ij}) \neq 0$ (which requires $N_{\mathrm{gen}}
\geq 3$), then $V_{ij} \neq V_{ij}^*$, and the forward phase
sequence differs from the $CP$-transformed sequence.  This phase
mismatch is $CP$ violation.
\end{proof}

\begin{proposition}[$CP$ Requires Three Generations]
\label{prop:three-gen}
$CP$ violation vanishes for $N_{\mathrm{gen}} < 3$.
\end{proposition}

\begin{proof}
For $N_{\mathrm{gen}} = 2$, the $2 \times 2$ CKM matrix has one
real parameter (the Cabibbo angle $\theta_C$) and no complex
phase.  $V_{ij} = V_{ij}^*$ for all $i, j$, so $CP$ is
preserved.

For $N_{\mathrm{gen}} = 3$, the $3 \times 3$ CKM matrix has
three angles and one irreducible phase $\delta$, giving
$\mathcal{J} = \mathrm{Im}(V_{us} V_{cb} V_{ub}^* V_{cs}^*)
\neq 0$.  (This is the Kobayashi--Maskawa theorem \cite{KM1973}.)
\end{proof}

\begin{remark}[Status of the RS derivation of $CP$ violation]
\label{rem:cp-status}
RS forces $N_{\mathrm{gen}} = 3$ from the face-pair count of the
$D = 3$ cube (three pairs of opposite faces).  Since the minimum
number of generations for a non-trivial CKM phase is $N_{\mathrm{gen}}
= 3$ (Kobayashi--Maskawa theorem), $CP$ violation is a \emph{structural
consequence} of the RS dimension-forcing.

However, this argument has two levels:
\begin{enumerate}[label=(\roman*)]
  \item RS derives $N_{\mathrm{gen}} = 3$ (via $D = 3$ cube geometry).
  \item The KM theorem (standard group theory) then implies that a
  $3 \times 3$ unitary matrix generically has one irreducible complex
  phase.
\end{enumerate}
Step (i) is the RS content; step (ii) is imported from the Standard Model.
The RS framework predicts the \emph{existence} of $CP$ violation (because
3 generations are forced and 3 is the minimum for KM phases), but does
not independently derive the \emph{magnitude} of the Jarlskog invariant
$\mathcal{J} \approx 3.2 \times 10^{-5}$.  Predicting $\mathcal{J}$
from the RS CKM matrix structure (rungs, gap weights, and three-generation
ladder structure) is an identified open problem.
\end{remark}

%% ============================================================
\section{$T$ Violation as the Complement}
\label{sec:t-violation}
%% ============================================================

\begin{corollary}[$CP$ Violation $\Leftrightarrow$ $T$ Violation]
\label{cor:cp-t}
Since $CPT$ is exact (Theorem~\ref{thm:cpt}), $CP$ violation implies
$T$ violation, and conversely.
\end{corollary}

\begin{proof}
If $CPT$ is a symmetry but $CP$ is not, then $T = (CP)^{-1}(CPT)$
cannot be a symmetry either (the composition of a non-symmetry with a
symmetry is not a symmetry).  Conversely, if $T$ is a symmetry and
$CPT$ is a symmetry, then $CP = (CPT) \cdot T^{-1}$ would also be a
symmetry, contradicting $CP$ violation.  Hence $CP$ non-invariance
$\Leftrightarrow$ $T$ non-invariance.
\end{proof}

\begin{remark}
The CPLEAR experiment~\cite{CPLEAR1998} and the BaBar/Belle
$B$-factories have directly observed $T$ violation in the neutral
kaon and $B$-meson systems, confirming this complementarity.
\end{remark}

%% ============================================================
\section{$CPT$ Tests}
\label{sec:tests}
%% ============================================================

\begin{center}
\renewcommand{\arraystretch}{1.3}
\begin{tabular}{lcc}
\toprule
\textbf{Test} & \textbf{Experimental bound} & \textbf{RS prediction} \\
\midrule
$|m_p - m_{\bar{p}}|/m_p$ & $< 7 \times 10^{-10}$~\cite{PDG2024} &
  0 (exact) \\
$|q_p + q_{\bar{p}}|/e$ & $< 7 \times 10^{-10}$~\cite{PDG2024} &
  0 (exact) \\
$|\tau_{\mu^+} - \tau_{\mu^-}|/\tau$ & $< 2 \times 10^{-5}$~\cite{PDG2024} &
  0 (exact) \\
$|m_{K^0} - m_{\bar{K}^0}|/m_{K^0}$ &
  $< 6 \times 10^{-19}$~\cite{PDG2024} & 0 (exact) \\
\bottomrule
\end{tabular}
\end{center}

%% ============================================================
\section{Predictions}
\label{sec:predictions}
%% ============================================================

\begin{enumerate}
  \item \textbf{$CPT$ exact}: All $CPT$ tests should yield null
  results.  Any $CPT$ violation would falsify RS.

  \item \textbf{$CP$ violation requires 3 generations}: Confirmed
  experimentally.

  \item \textbf{$T$ violation observed}: Confirmed by CPLEAR and
  $B$-factories.

  \item \textbf{No $CPT$-violating Lorentz violation}: SME
  (Standard Model Extension) coefficients for $CPT$ violation are
  all zero.
\end{enumerate}

%% ============================================================
\section{Conclusion}
\label{sec:conclusion}
%% ============================================================

$CPT$ invariance is automatic on the recognition ledger, following
from the individual cost-preserving nature of $C$, $P$, and $T$.
$CP$ violation arises because the 8-tick phase structure is not
preserved under $CP$ when weak-sector entries carry complex CKM
phases, which require $N_{\mathrm{gen}} \geq 3$.  Since RS forces
$N_{\mathrm{gen}} = 3$, $CP$ violation is a structural consequence:
three generations are necessary and sufficient for the KM mechanism
to produce an irreducible phase (see Remark~\ref{rem:cp-status}).

%% ============================================================
\begin{thebibliography}{99}

\bibitem{Christenson1964}
J.~H.~Christenson, J.~W.~Cronin, V.~L.~Fitch, and R.~Turlay,
``Evidence for the $2\pi$ Decay of the $K_2^0$ Meson,''
Phys.\ Rev.\ Lett.\ \textbf{13}, 138 (1964).

\bibitem{KM1973}
M.~Kobayashi and T.~Maskawa, ``CP-Violation in the
Renormalizable Theory of Weak Interaction,'' Prog.\ Theor.\
Phys.\ \textbf{49}, 652 (1973).

\bibitem{CPLEAR1998}
CPLEAR Collaboration, ``First Direct Observation of
Time-Reversal Non-Invariance in the Neutral-Kaon System,''
Phys.\ Lett.\ B \textbf{444}, 43 (1998).

\bibitem{RS_Axioms}
J.~Washburn, ``The Algebra of Reality: A Recognition Science Derivation
of Physical Law,'' \emph{Axioms} \textbf{15}(2), 90 (2026).
\texttt{doi:10.3390/axioms15020090}

\bibitem{PDG2024}
S.~Navas \textit{et al.}\ (Particle Data Group),
``Review of Particle Physics,'' Phys.\ Rev.\ D \textbf{110}, 030001
(2024).

\end{thebibliography}

\end{document}
